% ═══════════════════════════════════════════════════════════════════════════════
%  CAPÍTULO — Mediator
% ═══════════════════════════════════════════════════════════════════════════════

\chapter{Mediator}

\patroninfo{Comportamiento}{273}

% ─── Situación Problema ──────────────────────────────────────────────────────
\section{Situación Problema}

El sistema de mensajería de Airbnb conecta a huéspedes y anfitriones
dentro de la plataforma. Sin embargo, los mensajes no son simples textos:
pueden activar acciones automáticas. Cuando el anfitrión responde una
consulta de disponibilidad, el sistema debe verificar el calendario;
cuando un huésped confirma verbalmente el check-in, el sistema puede
actualizar el estado de la reserva; si alguna de las partes reporta
un problema, se abre automáticamente un ticket de soporte.

Si cada participante del chat —huésped, anfitrión, sistema de reservas,
soporte— se comunica directamente con los demás, se crea una red densa
de dependencias cruzadas. Añadir un nuevo tipo de participante
o una nueva acción automática requiere modificar todas las partes que
podrían interactuar con él.

% ─── Solución ────────────────────────────────────────────────────────────────
\section{Solución}

El patrón Mediator introduce \texttt{CentralMensajeria} como el único
punto de coordinación. Huéspedes, anfitriones y servicios automáticos
solo se comunican con el mediador, nunca directamente entre sí. El
mediador conoce las reglas de negocio: qué debe ocurrir cuando se
envía cierto tipo de mensaje y qué participantes deben ser notificados
o activados en consecuencia.

Esto reduce las dependencias entre participantes de $O(n^2)$ a $O(n)$,
hace explícitas y centralizadas las reglas de interacción, y permite
incorporar nuevos participantes o nuevas reglas de automatización
modificando únicamente el mediador.

% ─── Diagrama de clases ─────────────────────────────────────────────────────
\begin{figure}[H]
    \centering
    % \includegraphics[width=\textwidth]{imagenes/abstract_factory_diagrama.png}
    \caption{Diagrama de clases del patrón Abstract Factory}
\end{figure}


% ─── Roles de las Clases ────────────────────────────────────────────────────
\section{Roles de las Clases}

% Explicar la responsabilidad de cada clase/interfaz

\begin{itemize}
    \item \textbf{AbstractFactory:}
    \item \textbf{ConcreteFactory:}
    \item \textbf{AbstractProduct:}
    \item \textbf{ConcreteProduct:}
    \item \textbf{Client:}
\end{itemize}


% ─── Main ───────────────────────────────────────────────────────────────────
\section{Main}

% Explicar qué realiza el programa principal

\begin{lstlisting}[language=Java, caption=NombreClase]
 
\end{lstlisting}


% ─── Salida del Main ────────────────────────────────────────────────────────
\section{Salida del Main}

\begin{lstlisting}[language=Java, caption=NombreClase]
 
\end{lstlisting}


% ─── Clases ─────────────────────────────────────────────────────────────────
\section{Clases}

% ─── Clase 1 ────────────────────────────────────────────────────────────────
\subsection{NombreClase}

\begin{lstlisting}[language=Java, caption=NombreClase]
 
\end{lstlisting}


% ─── Clase 2 ────────────────────────────────────────────────────────────────
\subsection{NombreClase}

\begin{lstlisting}[language=Java, caption=NombreClase]
 
\end{lstlisting}


% ─── Clase 3 ────────────────────────────────────────────────────────────────
\subsection{NombreClase}

\begin{lstlisting}[language=Java, caption=NombreClase]
 
\end{lstlisting}


% ─── Conclusiones ───────────────────────────────────────────────────────────
\section{Conclusiones}

% Explicar:
% - Ventajas del patrón
% - Cuándo usarlo
% - Beneficios obtenidos
% - Posibles desventajas
% - Relación con principios SOLID